\documentclass[aspectratio=169]{beamer}
\usepackage{sudsley}

\title{System Design}
\subtitle{Code review and analysis}

% A code sample on the left, its review points revealed one by one on the right.
\lstset{basicstyle=\ttfamily\fontsize{6}{7.2}\selectfont}
\newenvironment{sample}[1]
  {\samplelabel{#1}\begin{columns}[T]\begin{column}{0.66\textwidth}}
  {\end{column}\end{columns}}
\newcommand{\review}{\end{column}\begin{column}{0.31\textwidth}\footnotesize}

\begin{document}

\maketitle

\begin{frame}{Recap: Homework 3}
  Complete the \texttt{ApiHandler} class in \texttt{src/handler/api.ts}, and submit:
  \begin{itemize}
    \item all source code, and
    \item a report (PDF) with a screenshot of the service's performance statistics and a
      short analysis.
  \end{itemize}
\end{frame}

\section{Engineering currency}

\begin{frame}{Engineering currency}
  \begin{enumerate}
    \item \textbf{Machine currency, the ``real'' cost.}
      CPU cycles, memory, database I/O, network calls.
      In industry, efficiency is money.
    \item \textbf{Human currency, the ``cognitive'' cost.}
      Readability, maintainability, time to understand.
      \begin{itemize}
        \item Code is read far more often than it is written.
        \item Cognitive load: how much effort does it take to see what the code does?
        \item Debug time: how fast can you fix it?
      \end{itemize}
  \end{enumerate}
\end{frame}

\begin{frame}{Machine currency}
  In the homework, efficiency is measured in \alert{units}. In the real world, units become
  a computing bill.

  \medskip
  An illustration. Suppose 1 unit = \$0.0001, over 10k requests:
  \begin{itemize}
    \item Implementation X costs 1.5 million units: \$150.00
    \item Implementation Y costs 140k units: \$14.00
  \end{itemize}
  \pause
  Scale to 1 million users (100$\times$ the requests), against a \$10k monthly budget:
  \begin{itemize}
    \item X's bill: \$15k a month. Over budget: \alert{bankruptcy.}
    \item Y's bill: \$1.4k a month. Well under: \alert{profit!}
  \end{itemize}
\end{frame}

\begin{frame}{What the simulation charges}
  Every request in the HW3 simulation also pays:
  \begin{itemize}
    \item 128 units to check the token with the Identity Provider
    \item 256 units for \emph{every} \texttt{getInstance()} call on the database table or an
      external client
  \end{itemize}
  So your totals will be well above the illustration. Compare them with each other.
\end{frame}

\begin{frame}{Human currency}
  Key metrics:
  \begin{itemize}
    \item Time to understand (TTU)
    \item Mean time to repair (MTTR)
  \end{itemize}
  Higher TTU means higher MTTR.
  \pause

  \medskip
  It's 3:00 AM and the production server is down.
  \begin{itemize}
    \item \textbf{Complex code:} two hours or more to fix, lost revenue, escalation to your
      manager.
    \item \textbf{Clean code:} under ten minutes to fix; customers may never notice.
  \end{itemize}
\end{frame}

\section{Code samples}

\begin{frame}[fragile]{Re-reading what you just wrote}
\begin{sample}{A}
\begin{lstlisting}
private handleRequestMachine(
    request: RequestMachineRequestModel,
): MachineResponseModel {
    const table = MachineStateTable.getInstance();
    const machines = table.listMachinesAtLocation(request.locationId);
    const available = machines.find(m => m.status === MachineStatus.AVAILABLE);
    if (!available) {
        return { statusCode: HttpResponseCode.NOT_FOUND, machine: undefined };
    }
    table.updateMachineJobId(available.machineId, request.jobId);
    table.updateMachineStatus(
        available.machineId, MachineStatus.AWAITING_DROPOFF);
    const updated = table.getMachine(available.machineId)!;
    this.cache.put(updated.machineId, updated);
    return { statusCode: HttpResponseCode.OK, machine: updated };
}
\end{lstlisting}
\review
\begin{itemize}
  \item<2-> Updates the database, then immediately queries it.
  \item<2-> An unnecessary read.
  \item<3-> \textbf{Why?} Lack of confidence in local state: the programmer wants to be
    ``sure'' they have the latest data.
\end{itemize}
\end{sample}
\end{frame}

\begin{frame}[fragile]{Building the result in memory}
\begin{sample}{B}
\begin{lstlisting}
private handleRequestMachine(
    request: RequestMachineRequestModel,
): MachineResponseModel {
    const table = MachineStateTable.getInstance();
    const machines = table.listMachinesAtLocation(request.locationId);
    var available = machines.find(m => m.status === MachineStatus.AVAILABLE);
    if (!available) {
        return { statusCode: HttpResponseCode.NOT_FOUND, machine: undefined };
    }
    const updatedMachine = {
        machineId: available.machineId,
        locationId: available.locationId,
        currentJobId: request.jobId,
        status: MachineStatus.AWAITING_DROPOFF
    }
    table.updateMachineStatus(updatedMachine.machineId, updatedMachine.status);
    table.updateMachineJobId(
        updatedMachine.machineId, updatedMachine.currentJobId);
    this.cache.put(updatedMachine.machineId, updatedMachine);
    return { statusCode: HttpResponseCode.OK, machine: updatedMachine };
}
\end{lstlisting}
\review
\begin{itemize}
  \item<2-> Constructs the updated item in memory.
  \item<2-> That costs nothing!
  \item<3-> \textbf{Takeaway:} if you own the data transformation, don't pay the database to
    repeat back what you just did.
  \item<4-> \textbf{Aside:} \texttt{var} makes readers check for reassignment. Prefer
    \texttt{const}.
\end{itemize}
\end{sample}
\end{frame}

\begin{frame}{A versus B}
  \centering
  \begin{tabular}{@{}lrr@{}}
    \toprule
    Operation & A (re-read) & B (in memory) \\
    \midrule
    List machines & 8 & 8 \\
    Update status & 15 & 15 \\
    Update job & 12 & 12 \\
    Get result & 8 (database read) & 0 (memory) \\
    Cache put & 6 & 6 \\
    \midrule
    \textbf{Total units} & \textbf{49} & \textbf{41} \\
    \bottomrule
  \end{tabular}

  \bigskip
  Database and cache work only; both also pay the same connection and token charges.
  Re-reading adds 8 units: about 20\% more than B.
\end{frame}

\begin{frame}{Two requests, one machine}
  A and B both read the list, pick an available machine, then write to it.
  \pause
  \begin{itemize}
    \item Two requests for the same location arrive together.
    \item Both pick the same machine. Both reserve it.
  \end{itemize}
  \pause
  \alert{How would you prevent that?} A conditional write that succeeds only if the machine is
  still \texttt{AVAILABLE} (compare-and-set), a lock, or a transaction.
\end{frame}

\begin{frame}[fragile]{Cache lookaside}
\begin{sample}{C}
\begin{lstlisting}
private handleGetMachine(
    request: GetMachineRequestModel,
): MachineResponseModel {
    const table = MachineStateTable.getInstance();
    const cached = this.cache.get(request.machineId);
    if (cached) return { statusCode: HttpResponseCode.OK, machine: cached };
    const doc = table.getMachine(request.machineId);
    if (!doc) return { statusCode: HttpResponseCode.NOT_FOUND };
    this.cache.put(request.machineId, doc);
    return { statusCode: HttpResponseCode.OK, machine: doc };
}
\end{lstlisting}
\review
\begin{itemize}
  \item<2-> \alert{This is good.}
  \item<3-> \textbf{Cache lookaside:} checks the cache first and skips the database on a hit.
  \item<4-> \textbf{Return early:} readable, no deep nesting.
  \item<5-> \textbf{But:} \texttt{getInstance()} runs before the cache check, so even a hit
    pays a 256-unit connection. Get the table only on a miss, or once in the constructor.
\end{itemize}
\end{sample}
\end{frame}

\begin{frame}{Is the cache worth it?}
  \centering
  {\small Counting only cache and database operations:}\par\smallskip
  \begin{tabular}{@{}lr@{}}
    \toprule
    Path & Units per call \\
    \midrule
    Cache hit & 0 \\
    Cache miss: miss 4 + database read 8 + cache write 6 & 18 \\
    No cache: database read & 8 \\
    \bottomrule
  \end{tabular}

  \bigskip
  \raggedright
  \pause
  Over 10k calls at a 50\% hit rate: \alert{90k units with the cache, 80k without.}

  \pause
  Caching pays off only above about a 56\% hit rate.
  What hit rate did \emph{your} simulation run report?
\end{frame}

\begin{frame}[fragile]{Extra work in the error path}
\begin{sample}{D}
\begin{lstlisting}
private handleStartMachine(
    request: StartMachineRequestModel,
): MachineResponseModel {
    // ... [Validation logic omitted] ...
    const smartMachine = SmartMachineClient.getInstance();
    try {
        smartMachine.startCycle(request.machineId);
        // ... [Success logic omitted] ...
    }
    catch (err) {
        db.updateMachineStatus(request.machineId, MachineStatus.ERROR);
        machine = db.getMachine(request.machineId);
        if (machine) {
            this.cache.put(request.machineId, machine);
        }
        return { statusCode: HttpResponseCode.HARDWARE_ERROR, machine };
    }
}
\end{lstlisting}
\review
\begin{itemize}
  \item<2-> On a hardware failure, control jumps to the \texttt{catch}.
  \item<2-> Writing \texttt{ERROR} is necessary; re-reading the machine is Sample A's extra
    read, on a path that is already failing.
  \item<3-> Bigger questions: is the machine left half-started? Should it call
    \texttt{forceStop}, retry, or time out?
\end{itemize}
\end{sample}
\end{frame}

\begin{frame}[fragile]{Idiomatic code that hides a cost}
\begin{sample}{E}
\begin{lstlisting}
private handleRequestMachine(
    request: RequestMachineRequestModel,
): MachineResponseModel {
    const db = MachineStateTable.getInstance();
    const availableMachines = db.listMachinesAtLocation(request.locationId)
        .filter(m => m.status === MachineStatus.AVAILABLE);
    if (availableMachines.length === 0) {
        return { statusCode: HttpResponseCode.NOT_FOUND };
    }
    const machine = availableMachines[0];
    // ... [Update logic omitted] ...
}
\end{lstlisting}
\review
\begin{itemize}
  \item<2-> What does \texttt{.filter(...)} do? Check the documentation (or ask an AI).
  \item<3-> It looks clean but hides an efficiency flaw: it walks the \emph{entire} array.
    A location with 1,000 machines checks all of them, every time.
  \item<4-> \texttt{.find(...)} stops at the first match. Same $O(N)$ worst case, but it exits
    early whenever an available machine comes first.
\end{itemize}
\end{sample}
\end{frame}

\begin{frame}{Takeaways}
  \begin{itemize}
    \item Once the logic works, check it briefly for efficiency.
    \item In an unfamiliar language or codebase, convenient-looking constructs are sometimes
      traps. Verify they do what you think.
    \item Ask (a search engine or an AI): is there a better way to do this?
  \end{itemize}
\end{frame}

\end{document}
